\section{Method: deriving requests from process state}
\label{sec:method}

\subsection{The planning object}
The source specification $K$ records process nodes and their conditions,
obligations, required facts, evidence capabilities and document routes. A
capability describes what the evidence must establish. A route is a subset
of the closed document catalogue $\mathcal D$ that can establish it; a
capability may have several alternative routes. Case state $s=(z,E)$
contains the assessment $z$ of the source conditions and the recorded
evidence $E$. The versioned planner returns documents, a next action and
the chains that justify them:
\begin{equation}
(D_t,a_t,J_t)=\Pi_K(s_t),\qquad s_{t+1}=U(s_t,e_{t+1}),
\label{eq:loop}
\end{equation}
Here $U$ admits an observation or correction and updates the state from
which the next plan is computed. The construction gives two guarantees.
An unresolved condition can produce a question but cannot trigger a
branch-specific document request (non-prematurity). Every document in $D_t$
comes from an active chain containing an obligation, a fact, a capability
and a route (provenance closure).

\subsection{Model judgements and controller decisions}
Model calls perform the semantic work. During preparation, fixed roles
extract propositions anchored to source quotations, synthesise the macro
process, refine propositions into guarded evidence rules and map
capabilities to the document catalogue. At inference time, the interpreter
judges whether each guard holds for the case. These roles return structured
judgements. Code activates obligations, projects requests, selects routes
and prioritises the next action, recording the justification for each
choice (Appendix~\ref{app:studya}). The model can interpret a passage or
assess a condition; document requests must pass through the controller.

\subsection{Study A: guarded rules over a closed catalogue}
\label{sec:v5}
We compiled the frozen Study~A specification from \nPassages{} exact
passages of Swiss insurance law and public household-insurance policies
and forms. The sources are Fedlex, CSS, AXA, Generali, Simpego and the
Swiss Insurance Association; the release records their URLs and content
hashes. Extraction produced \nPropositions{} operational propositions.
A deterministic check requires an evidence rule or an explicit
non-evidentiary classification for each proposition. The resulting
representation contains \nRules{} evidence rules, \nNonEvidentiary{}
non-evidentiary classifications, \nGuards{} case guards and \nCatalogue{}
document types. This completeness check covers the extracted inventory.

Each rule $r$ has a local applicability guard $g_r$, required capabilities
$\mathcal C_r$ and routes $\mathcal R_c$ for each capability $c$. For case
$x$, the interpreter assigns $z_g(x)\in\{\T,\F,\U\}$ to guard $g$.
A true or false verdict requires an exact quotation from that case;
otherwise the verdict becomes unresolved. Missing evidence also leaves a
guard unresolved. Each guard has its own model call, although case packets
are batched for execution, so individual branch decisions remain auditable.
Let $A(x)$ contain the unconditional rules and those whose guards are true.
Let $\mathcal R^{\rm conf}_g$ contain the mapped confirmation routes for a
guard that is true in the narrative but has not been established by admitted
evidence. Given held document types $H$, the planner returns
\begin{equation}
D(x,H)=\Bigl[\bigcup_{r\in A(x)}\bigcup_{c\in\mathcal C_r}\bigcup_{\rho\in\mathcal R_c}\rho
\;\cup\bigcup_{g\ \text{true, unconfirmed}}\ \bigcup_{\rho\in\mathcal R^{\rm conf}_g}\rho\Bigr]\setminus H .
\label{eq:v5}
\end{equation}
Thus a false guard adds no demand, an unresolved guard produces a question,
and the accepted verdicts fully determine the plan. Each request retains
its rule, fact, must-show statement and source spans.

Equation~\ref{eq:v5} also defines this version's limits. V5 unions the
mapped routes: it does not choose one sufficient alternative or assess a
held document's adequacy for a particular capability. Activation depends on
local rule guards. Changing macro-process edges while keeping the rules
and verdicts fixed leaves $D$ unchanged. In Study~A, an ``active path'' is
this case-conditioned applicability structure; it does not require
traversing every graph edge.

\subsection{Study B: applicability, acquisition and adequacy}
\label{sec:native-method}
Study~B deterministically compiles three public tenancy templates into a
representation shared by all arms. It uses no model-based source induction.
The case assessment distinguishes predicate truth, document presence,
individual adequacy and joint adequacy. Activation follows
\begin{equation}
q_s(z)=\phi_s(z)\land_3 J_s\bigl(\{q_p(z)\}_{p\in\mathrm{pa}(s)}\bigr),\ \ 
a_o(z)=q_{\mathrm{scope}(o)}(z)\land_3\phi_o(z),\ \ 
b_o(z)=a_o(z)\land_3\psi_o(z),
\label{eq:native-activation}
\end{equation}
where $J_s$ is the declared all/any join over parent scopes. Roots use the
empty all-join, which is true. The three-valued conjunction $\land_3$ is
false if any input is false, true if all are true and unresolved otherwise.
An obligation's applicability $a_o$ therefore depends on both its enclosing
scope and its own condition. The separate acquisition permission $\psi_o$
must also hold: only $b_o=\T$ permits an immediate request.

A route $\rho$ satisfies capability $c$ only when the required evidence is
present and adequate and the joint assessment is sufficient. Denote this
state by $\rho_c(E)$. Missing or inadequate documents form $M_{c\rho}$;
unknown presence or adequacy creates review items $V_{c\rho}$. The
controller selects
\begin{equation}
\rho_c^*=\operatorname*{arg\,min}^{\rm lex}_{\rho\in\mathcal R_c}
\bigl(\mathbf 1\{\rho_c(E)\ne\T\},\,|M_{c\rho}|,\,|V_{c\rho}|,\,
\operatorname{sort}(M_{c\rho}),\,\operatorname{id}(\rho)\bigr),
\label{eq:native-route}
\end{equation}
independently for each capability, then requests
$\bigcup_c M_{c\rho_c^*}$ over eligible obligations. This is a local route
choice, not a global minimisation. Inactive obligations produce no demand.
Unresolved eligibility produces conditional candidates and questions.
If no route exists, the controller records an evidence gap. Its action
priority is to resolve questions, review evidence, request a document,
review an active gap and, finally, choose an eligible substantive action.

This controller adds sufficiency semantics to the Study~A planning object.
It is a separate runtime, and the experiments do not compare its performance
with V5 or treat it as a repair of V5 (Appendix~\ref{app:correspondence}).
